\chapter{Related Work}
\label{chap:related_work}

The problem investigated in this thesis lies at the intersection of
population genomics, breed classification, genomic feature
selection, representation learning, and model interpretability.

This chapter reviews the main statistical and computational
approaches relevant to these areas.
Particular attention is given to methods that can produce
marker-level rankings, since their predictive utility is later
compared under a common experimental protocol in
Chapter~\ref{chap:experiments}.


\section{Population Genetics and SNP Genotyping}
\label{sec:pop_genetics_snps}

Single Nucleotide Polymorphisms (SNPs) are widely used genetic
markers for characterizing genomic variation within and between
populations
\cite{morin2004snps}.
High-throughput genotyping platforms enable tens of thousands of
loci to be measured simultaneously, providing genome-wide
representations of individual genetic variation in livestock
\cite{georges2019livestock, sorbolini2016signatures}.\
From a computational perspective, these data frequently produce a
$p \gg n$ setting in which the number of genomic variables is much
larger than the number of available individuals
\cite{saeys2007feature, chafai2023mlanimal}.
An additional property of genomic data is Linkage Disequilibrium
(LD), defined as the non-random association of alleles at different
loci
\cite{slatkin2008linkage, pritchard2001linkage}.LD creates correlation and redundancy among subsets of genomic
features.
This structure is relevant to dimensionality reduction and feature
selection because multiple correlated SNPs may carry partially
overlapping population information
\cite{slatkin2008linkage, saeys2007feature}.
Consequently, genomic marker reduction involves not only decreasing
the dimensionality of the input space, but also determining which
subset of markers retains useful discriminative information.


\section{Breed Classification and Genomic Selection in Livestock}
\label{sec:breed_classification_selection}

Genomic breed assignment aims to infer the population or breed of
origin of an individual from its molecular marker profile.

Such approaches have applications in population characterization,
conservation of livestock biodiversity, and genetic traceability
\cite{dimauro2013canonical, dalvit2007traceability,
      groeneveld2010diversity}.

The classification problem can be approached using both statistical
and machine-learning methods.
Previous studies have investigated multivariate discriminant
approaches and reduced SNP panels as mechanisms for distinguishing
livestock populations.

In cattle, Dimauro et al.\ used Canonical Discriminant Analysis
(CDA) and Discriminant Analysis (DA) to derive a reduced set of
highly discriminant SNPs for breed assignment and traceability
\cite{dimauro2013canonical}.
A closely related study extended the same general objective to sheep:
Dimauro et al.\ selected reduced SNP panels for discriminating
21 sheep populations and five geographic areas using the
OvineSNP50 BeadChip
\cite{dimauro2015sheep}.
These studies are particularly relevant to the present thesis because
they formulate marker reduction as a supervised population-
discrimination problem rather than as phenotype association.

A related but distinct field is genomic selection, which predicts
genetic merit or complex traits from genome-wide markers
\cite{meuwissen2001prediction}.
Although genomic selection is not the primary task addressed in this
thesis, methodological developments in this area are relevant because
they share the challenges of high dimensionality, correlated
markers, and limited population-specific sample sizes
\cite{georges2019livestock, chafai2023mlanimal}.

Multi-task variational approaches have also been proposed to learn
shared genomic representations across heterogeneous environments
and populations, although the VMGP framework used as architectural
inspiration in this thesis was originally validated in plant genomic
selection
\cite{zhao2025vmgp}.

The present thesis focuses specifically on breed classification,
sample efficiency, and marker reduction rather than on prediction
of continuous production traits.


\section{Statistical and Machine-Learning Approaches to SNP Ranking}
\label{sec:classical_marker_discovery}

Several approaches can be used to associate genomic markers with
population differentiation or predictive relevance.
The methods considered in this thesis span population-genetics
statistics, single-marker tests, multivariate discriminant
approaches, and machine-learning feature importance.


\subsection{Population Differentiation and Single-Marker Statistics}

The Fixation Index ($F_{ST}$) is a classical measure of genetic
differentiation among populations
\cite{wright1949genetical, weir1984estimating}.

When estimated at individual loci, $F_{ST}$ can be used to rank
markers according to differences in allele frequencies between
populations.

Single-marker association testing provides a related univariate
strategy.
Genome-Wide Association Studies traditionally evaluate the
association between genomic variants and phenotypic traits
\cite{hirschhorn2005genome}.

When large numbers of loci are tested independently, multiple
hypothesis testing becomes an important statistical consideration
\cite{storey2003significance}.

These methods provide interpretable marker-level statistics but
primarily evaluate each locus individually.
They therefore provide a useful baseline against which multivariate
and representation-learning approaches can be compared.



\subsection{Multivariate GWAS and Sample-Size Considerations}
\label{subsec:multivariate_gwas_sample_size}

The limitations of locus-wise genomic analysis under
high-dimensional and comparatively small-sample conditions have
also motivated explicitly multivariate formulations of GWAS.

Manca et al.\ proposed a multivariate GWAS (M-GWAS) procedure for
cattle based on the sequential use of Canonical Discriminant
Analysis, Discriminant Analysis, and Stepwise Discriminant Analysis
\cite{manca2020multivariate}.
The method was evaluated on both simulated and real data and, in the
simulation study, was compared with traditional single-marker and
Bayesian GWAS approaches over progressively increasing numbers of
animals.

Their results showed that the relative behavior of the methods
changed substantially with sample size: differences between
multivariate and traditional GWAS were larger in smaller-sample
settings and decreased as the number of animals increased
\cite{manca2020multivariate}.
This result should not be transferred directly to breed
classification, because the target problem in M-GWAS is
genotype--phenotype association rather than population assignment.
Nevertheless, it provides an important livestock-genomics precedent
for treating sample size as an explicit experimental variable and
for investigating methods that exploit joint marker structure.

The same work also situates discriminant approaches within a broader
line of livestock-genomics research that includes reduced marker
panels for cattle and sheep breed traceability and the detection of
selection signatures
\cite{dimauro2013canonical, dimauro2015sheep,
      sorbolini2016signatures}.


\subsection{Multivariate Discriminant Analysis and Selection Signatures}

Multivariate discriminant approaches model several variables jointly
and have been applied to genomic population analysis.

Canonical Discriminant Analysis (CDA), for example, constructs
canonical variables by maximizing separation among predefined
groups
\cite{sorbolini2016signatures}.

In livestock genomics, CDA-based analyses have been used both for
population discrimination and for deriving reduced marker panels.
Dimauro et al.\ reported that a reduced set of SNPs selected through
a discriminant framework retained sufficient information for
assignment of independent cattle samples
\cite{dimauro2013canonical}.

The same research line was subsequently applied to ovine data.
Using genotypes from 21 sheep populations, Dimauro et al.\ selected
compact panels for breed and geographic assignment and evaluated
their performance on validation data
\cite{dimauro2015sheep}.
This provides a particularly relevant classical baseline for the
ovine analyses developed in the present thesis.

Canonical loadings or correlations with canonical variables can
additionally be inspected to evaluate the contribution of individual
SNPs to discriminant dimensions
\cite{sorbolini2016signatures}.

These results motivate the broader idea that a high-density genomic
array may contain a considerably smaller subset of markers with
strong population-discriminative information.


\subsection{Machine-Learning Feature Importance}

Ensemble methods such as Random Forest (RF)
\cite{breiman2001random}
provide a non-linear alternative to locus-wise statistical ranking.

A Random Forest combines multiple decision trees trained on
different subsets of samples and features.
Feature relevance can subsequently be estimated using criteria such
as impurity-based importance or permutation importance
\cite{strobl2007bias}.

Different importance estimators can exhibit different statistical
properties and biases, particularly in the presence of correlated
predictors
\cite{strobl2007bias, nicodemus2010correlation}.

For this reason, Random Forest importance is treated in this thesis
as an empirical baseline rather than as an intrinsically weaker or
stronger alternative to gradient-based attribution.

The main statistical and machine-learning ranking paradigms considered
in the experimental comparison are summarized in
Figure~\ref{fig:classical_workflow_limitations}.
Although they all produce marker-level relevance information, they
derive it from different statistical or predictive criteria.

\input{figures/fig_classical_workflow_limitations.tex}





\section{Deep Learning Fundamentals: Autoencoders, VAEs, and Contrastive Learning}
\label{sec:deep_learning_fundamentals}

Deep representation-learning models provide an alternative to
operating directly in the original high-dimensional SNP space
\cite{bengio2013representation}.

Instead of evaluating each marker independently, these models learn
transformations of complete genotype vectors and can therefore
represent multivariate structure in the data.


\subsection{Generative Architectures and Variational Mapping}

Standard Autoencoders
\cite{rumelhart1986learning, hinton2006reducing}
consist of an encoder that compresses an input into a latent
representation and a decoder that attempts to reconstruct the
original observation.

Variational Autoencoders (VAEs)
\cite{kingma2013auto, rezende2014stochastic}
extend this formulation by introducing a probabilistic latent
variable $\mathbf{z}$ and by regularizing its distribution through
variational inference.

A relevant architectural reference for the generative model used in
this thesis is the Variational Multi-view Genomics Predictor (VMGP)
\cite{zhao2025vmgp}.

VMGP combines variational genomic representation learning with
multiple prediction objectives.
Its published evaluation concerns plant genomic selection; the
present work adapts the architecture to livestock population
classification and marker selection rather than assuming direct
domain transfer.


\subsection{Self-Supervised Contrastive Latent Spaces}

Contrastive representation learning follows a different objective.

Rather than reconstructing the input, the model learns
representations by controlling the similarity between related and
unrelated observations
\cite{oord2018representation, chen2020simple}.

In frameworks such as SimCLR, different augmented views of the same
sample form positive pairs, while other observations in the
mini-batch act as negatives
\cite{chen2020simple}.

The resulting representations are commonly normalized and compared
using cosine similarity.
Hyperspherical embedding spaces are particularly convenient in this
setting because normalized vectors lie on the unit sphere and their
dot product corresponds directly to cosine similarity
\cite{wang2020understanding}.

Applying these principles to genotype data requires domain-specific
choices regarding data augmentation and the definition of
biologically meaningful similarity.
Those choices are therefore specified explicitly as part of the
methodology in
Section~\ref{sec:spherical_contrastive_architecture}, rather than
being treated as established properties of generic contrastive
learning.



The exact architectures implemented for the experiments are
presented separately in
Sections~\ref{sec:vmgp_architecture} and
\ref{sec:spherical_contrastive_architecture}.



The conceptual difference between generative reconstruction and
contrastive organization of the latent space is summarized in
Figure~\ref{fig:generative_contrastive_topologies}.
The exact architectures implemented for the experiments are presented
separately in Sections~\ref{sec:vmgp_architecture} and
\ref{sec:spherical_contrastive_architecture}.

\begin{figure}[htbp]
\centering
\begin{tikzpicture}[
    >=Latex,
    font=\small,
    box/.style={draw, semithick, rounded corners=2pt, align=center,
                minimum height=10mm, minimum width=34mm, inner sep=4pt, fill=black!2},
    key/.style={box, fill=black!7},
    small/.style={draw, semithick, rounded corners=2pt, align=center,
                  minimum height=9mm, minimum width=27mm, inner sep=3pt, fill=black!2},
    arr/.style={->, semithick}
]
\node[font=\bfseries\normalsize] at (-4.1,3.25) {Generative};
\node[font=\bfseries\normalsize] at (4.1,3.25) {Contrastive};
\node[key] (gx) at (-4.1,2.35) {Input $\mathbf{x}$};
\node[box] (ge) at (-4.1,0.95) {Encoder};
\node[box] (gz) at (-4.1,-0.45) {Latent $\mathbf{z}$};
\node[box] (gd) at (-4.1,-1.85) {Decoder};
\node[key] (gr) at (-4.1,-3.25) {Reconstruction $\hat{\mathbf{x}}$};
\draw[arr] (gx)--(ge); \draw[arr] (ge)--(gz); \draw[arr] (gz)--(gd); \draw[arr] (gd)--(gr);

\node[key] (cx) at (4.1,2.35) {Input $\mathbf{x}$};
\node[small] (ca) at (2.35,0.85) {Augmented view A};
\node[small] (cb) at (5.85,0.85) {Augmented view B};
\node[small] (cea) at (2.35,-0.55) {Shared encoder};
\node[small] (ceb) at (5.85,-0.55) {Shared encoder};
\node[small] (cza) at (2.35,-1.95) {$\mathbf{z}_a$};
\node[small] (czb) at (5.85,-1.95) {$\mathbf{z}_b$};
\node[key] (closs) at (4.1,-3.45) {Contrastive objective};
\draw[arr] (cx)--(ca); \draw[arr] (cx)--(cb);
\draw[arr] (ca)--(cea); \draw[arr] (cb)--(ceb);
\draw[arr] (cea)--(cza); \draw[arr] (ceb)--(czb);
\draw[<->, semithick] (cza) -- node[midway, above=10pt, fill=white, inner sep=1pt]{positive pair} (czb);
\draw[arr] (cza)--(closs); \draw[arr] (czb)--(closs);
\draw[gray, dashed] (0,2.95)--(0,-4.05);
\end{tikzpicture}

\caption{Conceptual comparison between generative and contrastive representation learning. Generative models optimize reconstruction through a latent variable, whereas contrastive methods organize representations through similarity relationships between augmented views. Author's elaboration based on \cite{kingma2013auto,chen2020simple}.}
\label{fig:generative_contrastive_topologies}
\end{figure}



\section{Explainable AI: SHAP and Integrated Gradients}
\label{sec:explainable_ai_paradigms}

Deep models provide flexible non-linear representations but generally
do not expose the contribution of individual input features
directly
\cite{guidotti2018survey, arrieta2020explainable}.

Post-hoc feature-attribution approaches attempt to recover such
information after model training.

SHAP (Shapley Additive exPlanations)
\cite{lundberg2017unified}
provides an additive feature-attribution framework based on Shapley
values from cooperative game theory
\cite{shapley1953value}.

Exact Shapley-value evaluation is combinatorial in the number of
features, while practical SHAP implementations use approximation or
model-specific strategies to make attribution tractable
\cite{lundberg2017unified}.

Integrated Gradients
\cite{sundararajan2017axiomatic}
provides a gradient-based attribution method for differentiable
models.

Given an input $\mathbf{x}$, a baseline $\mathbf{x}'$, and model
output $F$, the attribution associated with feature $i$ is

\begin{equation}
\label{eq:integrated_gradients}
    \mathrm{IG}_i(\mathbf{x})
    =
    (x_i - x'_i)
    \int_{0}^{1}
    \frac{
        \partial
        F\left(
            \mathbf{x}'
            +
            \alpha(\mathbf{x}-\mathbf{x}')
        \right)
    }{
        \partial x_i
    }
    \,d\alpha .
\end{equation}

Integrated Gradients satisfies properties including sensitivity and
implementation invariance under the conditions defined in its
original formulation
\cite{sundararajan2017axiomatic}.

In genomic applications, sample-level attributions must subsequently
be aggregated into marker-level importance scores before a global
SNP ranking can be constructed.

The aggregation strategy adopted in this thesis is specified in
Chapter~\ref{chap:feature_importance} rather than assumed here.


\section{Research Gap}
\label{sec:research_gap}

The literature therefore provides established approaches for
population differentiation, genomic classification, multivariate
analysis, representation learning, and feature attribution
\cite{saeys2007feature, chafai2023mlanimal}.

However, these approaches are often developed for different
objectives or evaluated under different datasets and experimental
conditions.
Classical livestock studies already show that both marker selection
and the amount of available data can materially affect genomic
discrimination or association analyses
\cite{dimauro2013canonical, dimauro2015sheep,
      manca2020multivariate}.
Reviews of genomic machine learning likewise emphasize that there is
no universally superior modeling strategy and that data limitations,
standardization, and interpretability remain important challenges
\cite{chafai2023mlanimal, montesinoslopez2021deep}.

The present thesis focuses on a common experimental question:
how classification performance changes when both the amount of
training data and the dimensionality of the genomic representation
are progressively reduced.

Within the literature reviewed in this chapter, sample-size
analysis, marker reduction, representation learning, and
feature-attribution-based ranking are generally investigated under
different objectives or experimental settings.
The present work therefore evaluates these components within a single
controlled protocol that combines non-linear genomic representation
learning, progressive sample-size reduction, progressive marker-panel
reduction, and marker attribution.

Within this framework, different SNP-ranking strategies are compared
using the same downstream classifiers, data partitions, evaluation
metrics, and panel sizes.

This design makes it possible to distinguish three related but
different properties:

\begin{enumerate}

    \item the ability of a model to classify breeds using the
    complete genomic representation;

    \item the ability of a ranking method to identify compact SNP
    subsets that preserve predictive performance;

    \item the robustness of the observed behavior when training
    sample size is reduced or when the protocol is transferred to
    another livestock species.

\end{enumerate}

These considerations motivate the methodology introduced in
Chapter~\ref{chap:methods} and the experiments reported in
Chapter~\ref{chap:experiments}.
